\documentclass[12pt]{article}

\usepackage[utf8]{inputenc}
\usepackage[T1]{fontenc}
\usepackage[a4paper,landscape,margin=1.2cm]{geometry}
\usepackage{tikz}
\usepackage{tikz-uml}
\pagestyle{empty}

\begin{document}

\begin{tikzpicture}[scale=0.82, transform shape]

\node[font=\bfseries\large] at (-1.5, 8.2)
{Diagrama de classes - Decorator para itens do pedido};

\umlinterface[x=0,y=6]{ItemCardapio}{
}{
    + getNome() : String\\
    + getDescricao() : String\\
    + calcularValor() : double
}

\umlclass[x=-7.2,y=2.1]{Produto}{
    - idProduto : int\\
    - nome : String\\
    - ativo : boolean\\
    - descricaoProduto : DescricaoProduto
}{
    + getNome() : String\\
    + getDescricao() : String\\
    + calcularValor() : double
}

\umlclass[x=-12.5,y=2.1]{DescricaoProduto}{
    - idDescricao : int\\
    - descricaoDetalhada : String\\
    - valor : double
}{
    + getDescricaoDetalhada() : String\\
    + getValor() : double
}

\umlclass[type=abstract,x=5.2,y=2.1]{ItemCardapioDecorator}{
    \# itemDecorado : ItemCardapio
}{
    + ItemCardapioDecorator(itemDecorado : ItemCardapio)\\
    + getNome() : String\\
    + getDescricao() : String\\
    + calcularValor() : double
}

\umlclass[x=5.2,y=-2.9]{AdicionalDecorator}{
    - adicional : Adicional
}{
    + AdicionalDecorator(itemDecorado : ItemCardapio, adicional : Adicional)\\
    + getNome() : String\\
    + getDescricao() : String\\
    + calcularValor() : double
}

\umlclass[x=-1.1,y=-2.9]{Adicional}{
    - idAdicional : int\\
    - nome : String\\
    - valor : double
}{
    + getNome() : String\\
    + getValor() : double
}

\umlclass[x=-7.2,y=-2.9]{ItemPedido}{
    - quantidade : int\\
    - subtotal : double\\
    - produto : Produto\\
    - adicionais : List$<$Adicional$>$
}{
    + calcularSubtotal() : double\\
    - montarItemCardapio() : ItemCardapio
}

\umlimpl{Produto}{ItemCardapio}
\umlimpl{ItemCardapioDecorator}{ItemCardapio}
\umlinherit{AdicionalDecorator}{ItemCardapioDecorator}

\umlcompo[mult1=1,mult2=1,arg2=descricaoProduto,pos2=0.55]{Produto}{DescricaoProduto}
\umluniassoc[mult1=1,mult2=1,arg2=itemDecorado,pos2=0.65]{ItemCardapioDecorator}{ItemCardapio}
\umluniassoc[mult1=1,mult2=1,arg2=adicional,pos2=0.55]{AdicionalDecorator}{Adicional}

\umluniassoc[mult1=1,mult2=1,arg2=produto,pos2=0.65]{ItemPedido}{Produto}
\umlaggreg[mult1=1,mult2=0..*,arg2=adicionais,pos2=0.55]{ItemPedido}{Adicional}

\umlnote[x=-1.2,y=-6.7,width=15.5cm]{ItemPedido}{
Uso previsto: \texttt{ItemPedido} continua armazenando \texttt{Produto} e
\texttt{List\textless Adicional\textgreater}, preservando a estrutura atual do JSON. No
metodo \texttt{calcularSubtotal()}, ele monta a cadeia \texttt{ItemCardapio}:
\texttt{Produto} $\rightarrow$ \texttt{AdicionalDecorator} $\rightarrow$
outros adicionais. O subtotal passa a ser
\texttt{itemDecorado.calcularValor() * quantidade}.
}

\end{tikzpicture}

\end{document}
